\section{Local elementary factors and the least source genus}\label{sec:factors}

Let $S$ be the rational algebra of almost-complex cycle families
of \cite[\S2]{WueProduct}. Its relations include open excision and transport
of a compactly supported insertion over an interval. Multiplication
uses disjoint union of targets. Write $H$ for the Chern-number-zero,
virtual-dimension-zero part of Pardon's rational algebra of complex
one-cycles \cite{pardon}, and $\Phi:H\to S$ for the comparison.
We use $A_h$ for the balanced antidiagonal series of genus $h$ in $H$.
It has constant term one. The series $E_h=\Phi(A_h)$ is the
contribution of a neighborhood of an elementary curve of Ionel and
Parker \cite[\S3]{WueProduct}. Its coefficient of degree one is
$\Phi(x_{h,1})$, where $x_{h,1}$ is Pardon's genus-$h$ class of a
single curve of multiplicity one.

Let $C$ be a smooth embedded almost-complex curve of genus $g\ge2$
in a six-dimensional almost-complex germ with $c_1(TX)[C]=0$.
Assume a relatively compact tube around $C$ contains no other compact
irreducible curve and that the normal Cauchy--Riemann operator of $C$
is invertible on every connected finite \'etale cover, including $C$.
These conditions hold for the germs of Section~\ref{sec:spin}.
Choose a smaller disk tube. For $m>0$ let
$P_m$ be the class in $S$ of the characteristic function of the point
$m[C]$ in its degree-$m$ cycle space, and set $P=1+\sum_{m>0}P_mq^m$.
The variable $q$ labels $m[C]$ by $m$ under the retraction of the tube to $C$. An inclusion of two such
tubes preserves these classes by excision.
Write $h_d=1+d(g-1)$, and let $C_k(h)$ be the groupoid cardinality of
connected unramified degree-$k$ covers of a genus-$h$ curve. Each
isomorphism class contributes the reciprocal of its deck group order.
In particular $C_1(h)=1$.

\begin{proposition}\label{prop:factors}
For every such germ there are unique
integers $n_{d,h}$, of finite support in $h$ at each $d$, such that
\[
                  P=\prod_{d\geq1,\ h\geq0}E_h(q^d)^{n_{d,h}} .
\]
They vanish when $h<h_d$. If $c_m$ is the sum of the local determinant
signs divided by deck orders over the connected degree-$m$ covers,
then $n_d=n_{d,h_d}$ satisfies
\begin{equation}\label{eq:triangular}
           c_m=\sum_{d\mid m} n_d\,C_{m/d}(h_d).
\end{equation}
In particular this determines a full sequence of exponents of the
unique balanced product lifting $P$ to $H[[q]]$.
\end{proposition}

\begin{proof}
We first realize the neighborhood in a closed symplectic six-manifold
of vanishing first Chern class. In a product of two symplectic
two-tori take one horizontal torus and $g-1$ distinct vertical tori.
Smooth their $g-1$ positive transverse intersections. Locally the
smoothing of the coordinate axes is $zw=\delta$. In a fixed outer
annulus of each axis its graph may be cut off to the zero graph while
remaining symplectic, since the perturbation there is $C^1$-small
as $\delta\to0$. The annuli remove a disk from each incident torus
and glue the two boundary circles. The intersection graph is a tree,
so the resulting embedded symplectic surface has genus $g$. Regard it
as a surface in the product with a third two-torus, all three areas
chosen integral. Its class $A$ is primitive, since intersection with
a generic vertical torus in the first two factors pairs to one with
the four-dimensional surface class. Use the pullback of that pairing
to the six-manifold. In particular $d\mapsto dA$ is injective.

Identify the marked surface with this surface by an orientation-preserving
diffeomorphism. The symplectic normal bundle has first Chern number
$2g-2$. Choose a complex complement to $TC$ in the original tangent bundle along the curve and a Hermitian metric on that complement. A unitary normal identification with the model exists by the
classification of complex bundles over a surface. A tubular embedding
with this prescribed normal derivative carries the model's almost-complex
structure on the zero section to an almost-complex structure compatible
with the ambient symplectic form there. Any positive complex structure
on the oriented tangent plane is compatible with its area form.
Extend the compatible structure along the section to an ambient one
by extending a metric and using polar decomposition. On a small tube
the prescribed structure is tame and is uniformly close to the
extension, with equality along the section. Cut off their difference
in a chart of almost-complex structures to obtain a global tame
structure agreeing with the given germ on a smaller tube. This
construction uses the local almost-complex structure itself, including
its normal derivative.

Fix an energy bound. Take a smaller disk $V$ with closure in the
isolating tube. For every sufficiently close global tame structure,
there is no nonempty cycle contained in $\overline V$ whose support
meets $\partial V$ in this energy range. A violating sequence would
have a geometrically convergent subsequence by
\cite[Theorem~2.11]{diw}. Its limit at the prescribed structure
would be supported on $C$ and would meet $\partial V$, because
supports converge in the Hausdorff topology. Hence, in every positive
integral class, containment in $V$ selects a compact open and closed
insertion over a sufficiently short closed path to a generic tame
structure. At the first endpoint this is precisely the coefficient of
$P$. At the second endpoint the connected supports selected by $V$
form a union of finitely many clusters, after making the cluster
partition subordinate to this open and closed subset
\cite[Propositions~3.18--3.21]{diw}. All their classes are positive
multiples of $A$ since $V$ retracts onto $C$ and the symplectic area
is positive.

Apply the geometric excision and cluster recursion of
\cite[\S4]{WueProduct} to these selected clusters. This use of the
recursion is local: at the generic endpoint one includes the exterior
as an additional color to separate the supports, and retains the
factors of the selected colors. Each recursive correction is a
cluster in a proper multiple of its predecessor's class, by cluster
isotopy \cite[Theorem~3.25]{diw} and its cycle-space version in
\cite[\S4]{WueProduct}. Thus the remaining energy bound divided by the
core energy decreases at each recursion. Its leaves give an exact
finite product of elementary factors and inverses in $S$ at the
given cutoff. This part of the cluster argument uses a closed
symplectic ambient manifold and a bounded energy. The projective and
coherent-cohomology hypotheses of the compact complex product theorem
enter when one starts with the entire cycle series of a complex
target and lifts that identity to $H$.

At the first $q$-degree where two such products differ, the difference
is a sum of the degree-one classes $\Phi(x_{h,1})$ with coefficients
the differences of exponents. These classes are rationally independent
by the finite-genus test of \cite[Theorem~5.1]{WueProduct}.
Consequently the exponents stabilize as the cutoff increases,
are intrinsic to the germ, and the product of the $A_h(q^d)$ with
these exponents gives the stated lift. The uniqueness asserted here
is among balanced products with finite genus support at each degree.

Apply connected, fixed-target Gromov--Witten counts to the same
geometric decompositions, as in \cite[\S7]{WueProduct}. Compact open
localization, deformation of these insertions, and additivity on a
partition hold by \cite[Lemma~3.8 and Proposition~3.10]{diw}.
On disjoint targets the connected contributions add for normalized
factors, since the domain has an image in exactly one target and the
other insertion is evaluated on the empty cycle. These facts identify
the local connected series with the sum of the elementary connected
series with the very same integers $n_{d,h}$. This argument applies
to the given geometric relations on fixed targets. It requires no
enumerative pairing on arbitrary analytic parameter families.

For the original germ the count in degree $m$ vanishes below
$h_m$ by Riemann--Hurwitz, and at equality every stable map is a smooth unramified cover:
the arithmetic genus bound rules out multiple nonconstant components
and positive-genus ghosts, and a contracted rational tree would have
an unstable leaf. The tangent deformation complex of an unramified
map cancels against source deformations via $df$. The remaining
operator is the pulled-back normal operator, invertible by hypothesis.
Thus its determinant sign divided by the deck order gives the local
degree, and the sum is $c_m$. For a balanced elementary core of
genus $h\ge2$, the same argument gives vanishing below
$1+k(h-1)$ in degree $k$. At equality both spin lines coincide,
so every cover has sign $+1$ and the count is $C_k(h)$.
In degree one an elementary genus-$h$ factor starts at source genus
$h$ with coefficient one for every $h\ge0$
\cite[Proposition~3.3]{ip}. Now induct on $m$.
An elementary factor from $d<m$ enters degree $m$ only for $d\mid m$,
and by induction it has genus $h\ge h_d\ge2$. Thus its least possible
source genus in degree $m$ is
$1+(m/d)(h-1)\ge h_m$. Comparing the degree-one terms of the remaining
genus factors in order of increasing genus gives $n_{m,h}=0$ for
$h<h_m$. At equality the contributions from proper divisors occur
only for $h=h_d$. They have the values just computed, which proves
\eqref{eq:triangular}.
\end{proof}


\begin{corollary}\label{cor:cutoff}
For a homogeneous normal operator rigid through degree $M$, the same
construction gives the insertion of each pure multiple through $M$
and a unique balanced lift
\[
 P_T(D;q)=\prod_{j=1}^M\prod_s A_s(q^j)^{n^T_{j,s}(D)}
          \pmod {q^{M+1}} .
\]
Connected fixed-target integrals on these compact insertions
distinguish the truncated products.
\end{corollary}
\begin{proof}
A compact irreducible curve in the tube projects with positive degree,
since a fiber is a complex vector space. Its normalization gives a
kernel graph. Hence rigidity excludes all other curves through $M$,
which suffices for the finite-energy argument above. The connected
integrals respect open excision and continuation of the selected
compact clusters. In their product, a connected domain in a disjoint
union lies in one target. If two lifts first differ at label $j$, the
connected difference at that label is a finite sum
$\sum_s a_s(2\sin(v/2))^{2s-2}$. These functions have distinct
lowest powers, hence all $a_s$ vanish if the integrals agree.
\end{proof}
